\documentclass[10pt,a4paper]{amsart}
\usepackage[margin=19mm]{geometry}
\usepackage{amsmath,amssymb,mathtools}
\usepackage[scr=rsfs,cal=euler]{mathalfa}
\usepackage{xfrac}
\usepackage[unicode,hidelinks,bookmarks=false]{hyperref}
\newcommand{\ie}{i.e.\ }
\newcommand{\cf}{cf.\ }
\newcommand{\resp}{respectively\ }
\newcommand{\Subsection}{\subsection}
\providecommand{\nobreakdash}{\nobreak\hbox{-}}
\setlength{\emergencystretch}{2em}
\multlinegap0pt
\usepackage{siunitx}
\usepackage{siunitx}
\usepackage{siunitx}
\usepackage{siunitx}
\usepackage{amssymb}
\usepackage{float}
\usepackage{algorithm}\usepackage{algpseudocode}
\usepackage[shortlabels]{enumitem}
\usepackage{tikz}
\usepackage{siunitx}
\newtheorem{lemma}{Lemma}
\newcommand\bH{\mathbf{H}}
\newcommand{\bxi}{\boldsymbol{\xi}}
\newcommand{\bzeta}{\boldsymbol{\zeta}}
\newcommand\cN{\mathcal{N}}
\newcommand\vdiv{\mathop{\mathrm{div}}\nolimits}
\title{Fully mixed virtual element schemes for a new model of steady-state poroelastic stress-assisted diffusion in the brain}
\author{Isaac Bermudez, Bryan G\'omez-Vargas, Kent-Andr\'e Mardal, Andr\'es E. Rubiano, Ricardo Ruiz-Baier}
\date{Source: arXiv:2510.12307v3 (CC BY 4.0)}
\begin{document}
\maketitle

\noindent\emph{Source and licence.} Fully mixed virtual element schemes for a new model of steady-state poroelastic stress-assisted diffusion in the brain. Version arXiv:2510.12307v3 (CC BY 4.0), distributed by arXiv under the Creative Commons Attribution 4.0 International licence, http://creativecommons.org/licenses/by/4.0/, as recorded in the arXiv licence metadata for this exact version and shown on the abstract page https://arxiv.org/abs/2510.12307v3. Full licence notice: \texttt{licenses/arxiv-2510.12307-CC-BY-4.0.txt}.\par\smallskip

\noindent\textit{Editorial scope.} Counted unit: the article's lemma lem:boundedness (source0.tex lines 601--603). The article's complete original proof of it is delivered with it (source0.tex lines 604--620, one \texttt{proof} environment of 17 lines). No mathematics is written, repaired or re-derived: the statement and the proof are the article's own lines, in the article's own notation, and no retained line is substituted or re-flowed.  No further source-local context is needed: every cross-reference of every retained line resolves inside the two delivered spans.  Nothing was omitted from any retained span, so the package carries no omission record.  No retained line contains a citation, so the wrapper carries no bibliography and no bibliography entry.  No retained line contains a figure environment, an image-inclusion command or drawing code, so no image is delivered and the package declares no asset dependency.  Editorial formatting only, in addition: an AMS \texttt{amsart} preamble that restates the article's own declarations and macros used by the retained lines, namely the environments \texttt{lemma} on separate counters, together with the standard packages those lines need.  The retained environments print in this document as Lemma 1 in order of appearance, whereas the article prints the same items with the numbering of its own full document.  The complete native source of the article as distributed in the arXiv e-print for this exact version is bundled unchanged as \texttt{sources/187029/source0.tex}.
\begin{lemma}[boundedness]\label{lem:boundedness}
    The operator $\cN$ is bounded. That is, for any $M > 0$, there exists a constant $C_\star > 0$ such that $\|\cN(\bzeta)\|_{(\bH^4_{\mathrm{N}}(\vdiv,\Omega))'} \le C_\star$ whenever $\|\bzeta\|_{{{4,\vdiv;\Omega}}} \le M$.
\end{lemma}

\begin{proof}
    Let $M > 0$ and assume $\bzeta \in\bH^4_{\mathrm{N}}(\vdiv,\Omega)$ such that $\|\bzeta\|_{{4,\vdiv;\Omega}} \le M$. For any $\bxi \in \bH^4_{\mathrm{N}}(\vdiv,\Omega)$, we bound the duality pairing using the Cauchy-Schwarz and Hölder inequalities
    \begin{align*}
        |[\cN(\bzeta), \bxi]|
        \le \rho_2 \|\bzeta\|_{0, \Omega} \|\bxi\|_{0, \Omega} + \eta \|\bzeta\|_{0,4;\Omega}^3 \|\bxi\|_{0,4;\Omega} + \|\vdiv \bzeta\|_{0,\Omega} \|\vdiv \bxi\|_{0,\Omega}.
    \end{align*}
    Next, thanks to the continuous Sobolev embedding $\mathbf{L}^4(\Omega) \hookrightarrow \mathbf{L}^2(\Omega)$, there exists a constant $C_{\text{emb}} > 0$ such that
    \begin{align*}
        |[\cN(\bzeta), \bxi]| &\le \rho_2 C_{\text{emb}}^2 \|\bzeta\|_{0,4;\Omega} \|\bxi\|_{0,4;\Omega} + \eta \|\bzeta\|_{0,4;\Omega}^3 \|\bxi\|_{0,4;\Omega} + \|\vdiv \bzeta\|_{0,\Omega} \|\vdiv \bxi\|_{0,\Omega} \\
        &\le \left( \rho_2 C_{\text{emb}}^2 \|\bzeta\|_{{4,\vdiv;\Omega}} + \eta \|\bzeta\|_{{4,\vdiv;\Omega}}^3 + \|\bzeta\|_{{4,\vdiv;\Omega}} \right) \|\bxi\|_{{4,\vdiv;\Omega}},
    \end{align*}
    which says that
    \begin{equation*}
        \|\cN(\bzeta)\|_{\bH^4_{\mathrm{N}}(\vdiv,\Omega)'} \le (\rho_2 C_{\text{emb}}^2 + 1) M + \eta M^3 := C_\star.
    \end{equation*}
    Since $C_\star$ is finite and depends only on $M$ and physical parameters, $\cN$ maps bounded sets to bounded sets.
\end{proof}

\end{document}
